%------------%
%->> Section 4.5: 本章小结
%--------%

\section{本章小结}

本章围绕故障诊断智能体平台的设计与实现，说明了第三章所提出方法在系统中的落地方式。首先，本文给出了平台的总体架构，说明了前端展示层、Agent 执行层、Skill 组织层、MCP 接入层、后端服务层和数据存储层之间的分工关系，并解释了以 OpenCode 作为执行核心的原因。其次，本文分别说明了故障指纹构建、SOP 自动化构建和在线故障诊断三类方法在平台中的实现位置及其对应的数据流。随后，围绕工作区、任务实体和知识对象，给出了在线诊断任务与离线知识构建任务的执行过程和数据管理方式。最后，本文介绍了诊断执行界面、工作区界面和故障指纹界面，说明了平台如何支持诊断过程展示、案例级任务管理和故障指纹知识查看。

通过上述设计，第三章提出的故障指纹检索、SOP 增量更新和在线诊断执行被组织到同一平台之中，并通过工作区和结果写回机制形成闭环。第五章将在此基础上开展实验验证与分析。